\section{Le Web de données et les nouveaux protocoles}
Avant la diffusion des données était davantage passive où les bibliothèques numériques interrogeait les entrepôts OAI afin de moisssonner leurs idées, sans forcément intéragir directement avec celle-ci. La transition de la bibliothèque numérique de la BIU Cujas vers le Web sémantique ne saurait être réduite à une simple mise à jour d'infrastructure : elle engage une véritable refondation de l'objet documentaire lui-même. Aujourd'hui, le document numérique rompt définitivement avec la logique de faire la numérisation sous forme de captures d'image. Désormais, la représentation du document numérisé est devenu plus complexe. Comme l'a théorisé Jean-Michel Salaün au sein du collectif Roger T. Pédauque, \enquote{il est probable que le document numérique donne une dimension nouvelle à la notion d’intertextualité. Celle-ci se prolonge dans l’inter-sémioticité. L’inter-sémioticité du document numérique qui intègre texte, écrit, image et son ne change pas la nature du texte, elle en accroît en revanche la complexité. Les rapports d’inclusion réciproque constituent un objet de recherche nouveau qui devrait mieux spécifier la relation entre sémantique et sémiotique}. Le document numérique est devenu un objet structurée en trois dimension : à travers la forme, le texte et l'interconnexion sémantique. C'est pourquoi, faire entrer le patrimoine juridique dans le Web de données (Linked Open Data) consiste à matérialiser cette déconstruction pour substituer au document-objet traditionnel un objet dynamique et ouvert d'entités liées.

\subsubsection{La rigueur de la structure logique}
La structure logique des documents numérisés constitue le fondement de toute bibliothèque numérique patrimoniale. Comme le rappelle Isabelle Westeel, à partir des années 1970, l'émergence historique des langages de balises vient \enquote{de la nécessité de structurer l'information}, en se préoccupant \enquote{du contenu du document et non pas de sa mise en forme (contrairement au HTML qui gère des balises de mise en forme)}. Cette évolution consacre une distinction majeure dans la composition du document entre \enquote{une structure logique (contenu organisé) et d'une structure physique (mise en forme de l'information)}. Toutefois, la simple conformité syntaxique n'est pas suffisante pour garantir l'interopérabilité des données. Le document est \enquote{bien formé}, lorsqu'il respecte \enquote{l'ensemble des règles de la recommandation} ou l'utilisation de la DTD \textit{Document Type Definition} et des schémas XML. La DTD liste les éléments, les attributs associés dans le document, participant l'interopérabilité des métadonnées.

Ce sont ces mêmes métadonnées qui sont moissonnées à travers les entrepôts d'OAI-PMH des bibliothèques. Actuellement, l'entrepôt d'OAI-PMH de CujasNum rencontre des problèmes de moissonnages. L'une des raisons vient de l'absence de balises manquantes obligatoires comme \enquote{earliestDatestamp}, bloquant le moissonnage automatique par Gallica. Face à ce problème récurrant auquel est confronté la bibliothèque Cujas, l'établisssement souhaite qu'avec la nouvelle bibliothèque le moissonnage se fera de nouveau. En effet, à cause de l'impossibilité de faire moissonner leurs collections par les institutions nationales, comme la BNF, cela entraîne un manque de visibilité de leurs documents dans le Web de données.

\subsection{Le dépassement de l'OAI-PMH : de la diffusion passive aux données liées}
Depuis les années 2000, le protocole OAI-PMH (\textit{Open Archives Initiative Protocol for Metadata Harvesting})couplé au format Dublin Core simple a constitué la brique standard incontournable pour le moissonnage des métadonnées sur le Web. Si cette simplification transitoire a permis une diffusion à bas coût de l'information bibliographique, elle se heurte aujourd'hui aux exigences scientifiques des bibliothèques de recherche patrimoniale. Comme le démontre Vincent de Lavenne de la Montoise, le Dublin Core simple, limité à quinze éléments, est incapable de reproduire la richesse des documents complexes et entraîne \textit{de facto} une disparition des liens d'autorité. En effet, le Dublin Core simple ne possède \enquote{aucun mécanisme permettant de représenter ces notices d'autorité qui concerne des entités non documentaires (personnes, concepts, etc.)}. Les notices d'autorité sont la base permettant de créer des liens entre les notices bibliographiques. A la bibliothèque Cujas, cette \enquote{mise à plat} imposée par ce format provoque une destruction de l'arborescence initiale des documents complexes. Un périodique de doctrine, par exemple, s'aritcule obligatoirement selon une arborescence stricte  (titre de la revue, année de publication, fascicule, article) que le système actuel peine à restituer de manière interopérable, empêchant ainsi la bibliothèque de mettre en avant ce fonds numérisés.

Vincent de Lavelle rappelle que ce format a été adopté massivement, environ \enquote{74 \% des forunisseurs ne dépassaient pas l'utilisation de cinq éléments sur les quinze qui sont dc: creator, dc: identifier, dc: title, dc: date et dc: type}, participant à \enquote{l'exclusion d'enrichissement}. La souplesse qu'offre le Dublin Core inquiète les fournisseurs de données face à cette perte de l'information, amenant à \enquote{concacténer des données diverses dans des champs textuels}. La balise <dc: publisher> en est le parfait exemple, car il est souvent arriver au bibliothèque d'y concaténer le nom de l'éditeur et le lieu de la publication. Il faut donc envisager d'autres moyens de décrire les documents sur le Web en ayant, par exemple, recours au format de sérialisation.

\subsubsection{Décrire les documents sur le Web}
Si le protocole OAI-PMH a permis de diffuser massivement des métadonnées sous forme de notices plates, la transition vers le Web sémantique impose de repenser la description même des ressources. Décrire les documents sur le Web ne consiste plus à rédiger une notice bibliographique isolée, mais à insérer chaque entité (personne, concept) dans un réseau mondial de données sémantiques. 

Cette transformation s'appuie sur le modèle RDF (\textit{Resource Description Framework}), qui s'affranchit de la rigidité des bases de données traditionnelles (SQL ou XML) pour modéliser l'information d'un graphe décentralisé composé de triplets (Sujet, Prédicat, Objet). Dans cette architecture sémantique, la description, par exemple, d'un traité de la listet Pfister-Roumy devient malléable et évolutive: l'ajout de nouvelles propriétés ou de liens vers des référentiels externes s'effectue par simple incrémentation de triplets, sans risque de briser la structure logique du système d'information. Cela permettrait de mieux s'adapter aux pratiques des chercheurs, en termes de données. En 2017, les spécialistes d'ingénierie documentaire s'interrogeait sur le format de données à proposer aux chercheurs ainsi que sur leurs pratiques. Par ailleurs, la bibliothèque Cujas s'interroge également dessus souhaitant davantage proposer un format de données correspondant aux pratiques de recherche des chercheurs. Antoine Courtin plaidait pour un grand réalisme pragmatique en rappellant que \enquote{les pratiques des chercheurs, en termes de données, se limitent bien souvent à copier des données dans un document de traitement de texte. Dans ce contexte, est-il opportun de chercher à adopter des formats, des langages de requête, en somme des logiques qui ne sont pas ceux des communautés que l'on sert?}. Finalement, au lieu d'imposer des formats sémantiques trop complexes et rigides aux chercheurs, les bibliothèques doivent s'adpater à leurs gestes quotidiens les plus simples, comme le copier-coller de texte ou de référence. Ils cherchent avant tout à manipuler des données brutes. 

L'usage des technologies du Web sémantique a un potentiel fort d'amélioration des services rendus aux usagers. Ces technologies ouvrent la voie à l'adoption d'autres standards notamment le protocole IIIF permettant la multi-diffusion des images. 

\subsection{Le protocole IIIF et la multi-diffusion}
Grâce au protocole IIIF (\textit{International Image Interoperability Framework}), la bibliothèque Cujas pourrait reconstruirre virtuellement ses collections, comme la liste Pfister-Roumy, en appellant les manifestes d'images stockés à la BNF, offrant une alternative scientifique robuste aux \enquote{Essentiels du Droit} de Gallica sans héberger physiquement les fichiers. En effet, c'est l'une des ambitions de l'établissement, à travers cette refonte, elle envisage de devenir une alternative aux documents de droit exposés sur Gallica. L'attente interne n'est plus de se contenter d'exposer passivement ses notices bibliographiques sur sa bibliothhèque numérique, mais d'interconnecter activement ses notices dans le système documentaire. En adoptant ce protocole pour sa nouvelle bibliothèque numérique, la bibliothèque Cujas transforme alors le rôle de l'usager. Il n'est plus passif face à l'information transmise, mais devient un acteur de celle-ci en lui permettant à travers les outils misent à sa disposition de s'en servir pour contribuer à la valorisation des documents numérisés. Aujourd'hui, \enquote{les outils du web 2.0 ouvrent ainsi une brèche dans l'implication des utilisateurs dans la vie de la bibliothèque numérique, en leur permettant de contribuer à la valorisation et à la diffusion des collections}. L'objectif serait alors d'offrir aux utilisateurs la possibilité de contribuer à des tâches de traitement des documents, en leur permettant de faire de la transcription, ou de corriger des textes océrisés ou d'annoter les documents.

Pour répondre à ce besoin d'ouverture et de mutualisation, l'ingénierie documentaire recommande de se tourner vers les technologies du web sémantique et du protocole IIIF. L'utilisation du protocole IIIF permettrait d'alléger considérablement la charge des serveurs de la bibliothèque tout en permettant d'afficher, dans une même visionneuse experte, un volume numérise par la bibliothèque Cujas à côté d'un volume numérisé par la BNF. En utilisant le protocole IIIF, il permet de \enquote{manipuler les documents de manière décentralités}, ainsi grâce aux manifestes l'usager n'est plus \enquote{un récepteur passif}. En effet, grâce à cette outil, il est possible \enquote{d'annoter les images, comparer des feuilles à distances, et concevoir des parcours de recherche}. Le passage massif au protocole IIIF et à l'utilisation de la visionneuse Mirador ou Universal Viewer, s'observe chez la majorité des bibliothèques numériques créées aujourd'hui par les institutions patrimoniales. 

La transition vers le Web de données permettrait à la bibliothèque Cujas de s'inscrire dans un noeud dynamique au sein d'une communauté de bibliothèques entraînant une meilleure visibilité du contenu numérisé par l'établissement sur le Web, favorisant ainsi les opérations de recherche des usagers. La transition bibliographique est la traduction nationale de cette ambition, entraînant une réflexion nationale autour des codes de catalogage, comme RDA-FR, et à son éventuelle adoption.